% !PRINTED 402
\weberpage{402}
%<<CONT

4. \gesperrt{If $\mathfrak{a}$ is an ideal coprime to $Q$, then every pair of
numbers $x_1, x_2$ satisfying the conditions (12) furnishes us a positive
number $y$ divisible by $\mathfrak{a}$ with positive norm of the order $[Q]$}

\weq{13}{y = a\,x_1 + \frac{b+\sqrt{D}}{2}x_2,}

\gesperrt{and among these numbers no two are associated.}

\begin{center}
\S\,111. \gesperrt{The Dirichlet Limit Formula.}
\end{center}

If $t$ denotes a positive constant, then by the condition:

\weq{1}{N(y) = a\,(a\,x_1^2 + b\,x_1 x_2 + c\,x_2^2) \lessgtr t,} % Source prints a combined less/greater relation; \lessgtr is the matching corpus macro.
and, when the discriminant is positive, by the isolation conditions, there is
bounded, in a plane in which $x_1, x_2$ are rectangular coordinates, a region
$F_t$ which for negative discriminant is bounded by an ellipse, and for
positive discriminant by an arc of a hyperbola and by straight lines. The
points whose coordinates $x_1, x_2$ are integral rational numbers are called \gesperrt{lattice points}.

To a whole system of associated numbers whose norm is positive and less than
$t$ there then corresponds, for positive discriminant, one lattice point in
$F_t$, and for negative discriminant two, and in the two exceptional cases
four and six lattice points.

If $Z_t$ denotes the number of lattice points in the region $F_t$, and $V$ the
area of the region $F_t$, then [Bd. II, § 194, (6)]:

\weq{2}{V = Z_t\,t^{-1} + R_t\,t^{-\frac{1}{2}},}
in which $R_t$ does not become infinite as $t$ increases infinitely.

For negative discriminant one obtains $V$ from the well-known area of the
ellipse:

\weq{3}{V = \frac{2\pi}{a\sqrt{-D}},}
and for positive discriminant one obtains it most simply by introducing,
according to § 110, (7), in place of $x_1, x_2$ the variables of integration
$\xi_1, \xi_2$, whose limits are $0, 1$ and $-\infty, 0$:

\weq{4}{V = \frac{\log\varepsilon}{a\sqrt{D}}.}
